%% November 4, 2006 -- Updated for Version 1.1 installation
\documentclass[12pt,a4paper]{article}

\usepackage[english]{babel}
\usepackage[latin1]{inputenc}
\usepackage[T1]{fontenc}
\usepackage{listings}
\usepackage{epic,eepic}

\setlength{\oddsidemargin}{0cm}
\parindent = 0pt
\parskip = 1pc

\newcommand{\memberitem}[1]{
                              \vspace{-0.5pc}
                              \end{minipage}\\
                              \phantom{M}#1:\\
                              \phantom{MM}
                              \begin{minipage}[t]{12cm}

                           }
\newcommand{\memdsc}{\memberitem{Description}}
\newcommand{\meminp}{\memberitem{Parameters}}
\newcommand{\memdef}{\memberitem{Constructor settings}}
\newcommand{\memret}{\memberitem{Return values}}

\newenvironment{MemberDescr}
  {
    \begin{minipage}{15cm}
    \bf
  }
  {
    \end{minipage}
  }

\newcommand{\pntr}{\raisebox{-0.25em}{*}\ }
\newcommand{\ppntr}{\raisebox{-0.25em}{**}\ }


\begin{document}

  \lstset{language=C++,
          extendedchars=true,
          showstringspaces=false,
          stringstyle = \ttfamily,
          basicstyle=\footnotesize
         }

  \thispagestyle{empty}

  \title{\Huge \textbf{Extension to ePiX}}
  \author{\\ \\ Svend Daugaard Pedersen}
  \maketitle

  \newpage
  \tableofcontents
  \newpage

  \pagestyle{headings}  % empty plain headings myheadings

  \section{Introduction}

  ePiX\_ext is a set of extensions to the ePiX library, offering a
  Cartesian and logarithmic coordinate systems (both single and
  double) with a lot of parameters that can be set by the
  user. Further more it offers drawing of a hatched polygon as well as
  hatched area between two function graphs.

  You need some familiarity with ePiX to be able to use this extension.

  \section{Installation}

  To install ePiX\_ext, use the option \verb+--with-contrib+ when
  you configure the main package.  The library will be compiled
  separately (\texttt{libepixext.a}), and the library and header will
  be installed automatically when you do \texttt{make install}.

  To use these extensions in a source file, the header must be
  included explicitly in each source file's preamble:
  \begin{verbatim}
    #include "epix.h"
    #include "epix_ext.h"

    using namespace ePiX;
    using namespace ePiX_contrib;
  \end{verbatim}

  \section{Cartesian Coordinate System}

  The declaration statement

  \begin{lstlisting}
    CartesianCoord cs(-3,5,-2,6);
  \end{lstlisting}

  creates a cartesian coordinate system with x-bounds $-3\ldots5$ and
  y-bounds $-2\ldots6$. To draw the coordinate system using the
  standard layout set by the constructor of \textit{CartesianCoord},
  execute the member function \textit{draw} in \textit{cs}:

  \begin{lstlisting}
    cs.draw();
  \end{lstlisting}

  Here is an example of a complete C++ program making such a standard
  coordinate system:

  \begin{minipage}{7.5cm}
    \begin{lstlisting}
  #include "epix.h"

  int main()
  {
    CartesianCoord cs(-3,5,-2,6);

    picture(P(12, 12));
    unitlength("5mm");

    begin();
    cs.draw();
    end();

    return 0;
  }
    \end{lstlisting}
  \end{minipage}
  \begin{minipage}{5.5cm}
    \input{graf01.eepic}
  \end{minipage}

  To make more division points you have to execute the member function
  \textit{xMark()} and/or \textit{yMark()} with the position of the first
  mark and the number of marks as parameters.

  \begin{minipage}{7.5cm}
    \begin{lstlisting}
  #include "epix.h"

  int main()
  {
    CartesianCoord cs(-3,5,-2,6);

    picture(P(12, 12));
    unitlength("5mm");

    begin();
    cs.xMarks(-2,6);
    cs.yMarks(-1,6);
    cs.draw();
    end();

    return 0;
  }
    \end{lstlisting}
  \end{minipage}
  \begin{minipage}{5.5cm}
    %\begin{center}
      \input{graf02.eepic}
    %\end{center}
  \end{minipage}

  The member functions \textit{xMark()} and \textit{yMarks()} accept a
  third parameter spe\-ci\-fying the the distance between the
  marks. The default value of this is 1.

  The member functions \textit{xLabels()} and \textit{yLabels()} are
  setting other labels on the axis than the default. These member
  functions have one, two or three parameters. The first (mandatory)
  parameter is a NULL ter\-mi\-na\-ted array of string pointers. The
  second parameter is used to specify the position af the first label
  (default is position of the first mark) and the third the distance
  between the labels (default is distance between marks).

  \begin{minipage}[c]{7.5cm}
    \begin{lstlisting}
  #include "epix.h"

  static char* xlabels[] = {"-2","-1","","1","2","3",NULL};
  static char* ylabels[] = {"-1","","1","2","3","4",NULL};

  int main()
  {
    CartesianCoord cs(-3,5,-2,6);

    picture(P(12,12));
    unitlength("5mm");

    begin();
    cs.xMarks(-2,6);
    cs.yMarks(-1,6);
    cs.xLabels(xlabels);
    cs.yLabels(ylabels);
    cs.draw();
    end();

    return 0;
  }
    \end{lstlisting}
  \end{minipage}
  \begin{minipage}{5.5cm}
    \vspace{5pc}
      \input{graf03.eepic}
  \end{minipage}

  Note the empty string at the 0-position to avoid collision between a label
  and the axis.

  The length and position of the marks as well as the position of the
  labels can be specified. The position values are POSITIVE, NEGATIVE
  and (for the marks) CENTER. The marks are placed with the member
  function \textit{markLayout()} with four parameters: the mark length
  (in pt) and the position for x-axis and the same for the y-axis. The
  labels position is set using the member function \textit{labelPos()}
  with one or two parameters. The one parameter version sets the
  position to the same for both axis.

  \begin{minipage}{7.5cm}
    \begin{lstlisting}
   #include "epix.h"
   int main()
   {
     CartesianCoord cs(-3,5,-2,6);
     picture(P(12,12));
     unitlength("5mm");
     begin();
     cs.markLayout(2,POSITIVE,2,NEGATIVE);
     cs.labelPos(NEGATIVE);
     cs.draw();
     end();
     return 0;
   }
    \end{lstlisting}
  \end{minipage}
  \begin{minipage}{5.5cm}
    \input{graf04.eepic}
  \end{minipage}

  The axis name and the position of the name can be set with the member
  functions \textit{xName()} and \textit{yName()}. The parameters are the
  name and the position.% (use \textbf{NULL} for the name to keep the default).

  \begin{minipage}{7.5cm}
    \begin{lstlisting}
  #include "epix.h"
  int main()
  {
    CartesianCoord cs(-3,5,-2,6);
    picture(P(12,12));
    unitlength("5mm");
    begin();

    cs.xName("t/s",POSITIVE);
    cs.yName("E/kJ",NEGATIVE);
    cs.draw();

    end();
    return 0;
  }
    \end{lstlisting}
  \end{minipage}
  \begin{minipage}{5.5cm}
    \input{graf05.eepic}
  \end{minipage}

  The CENTER position of the names places the names at the end of the
  axis. If this position is used you probably need to shorten the
  length of the axis (or the names will be written outside the picture
  bounds). This is done by calling the member function
  \textit{axisBounds())}. This function accepts one or two
  parameters. The one-parameter version will set the upper right
  "corner" (most probably what you need), or both lower left and upper
  right corner.

  \begin{minipage}{7.5cm}
    \begin{lstlisting}
  #include "epix.h"
  int main()
  {
    CartesianCoord cs(-3,5,-2,6);
    picture(P(12,12));
    unitlength("5mm");
    begin();

    cs.axisBounds(P(4,5));
    cs.xName(CENTER);
    cs.yName(CENTER);
    cs.draw();

    end();
    return 0;
  }
    \end{lstlisting}
  \end{minipage}
  \begin{minipage}{5.5cm}
    \input{graf06.eepic}
  \end{minipage}

  Also the crossing point for the axis can be specified (default is
  (0,0)) and you can ask for a "broken" axis:

  \begin{minipage}{7.5cm}
    \begin{lstlisting}
  #include "epix.h"
  static char* ylabels[] = {"22","23","24","25","26",NULL};
  int main()
  {
    CartesianCoord cs(-1,7,19,27);
    picture(P(12, 12));
    unitlength("5mm");
    begin();

    cs.axisCross(P(0,20));
    cs.xMarks(1,6);
    cs.yMarks(22,5);
    cs.yLabels(ylabels);
    cs.yName(POSITIVE);
    cs.yBroken();
    cs.draw();

    end();
    return 0;
  }
    \end{lstlisting}
  \end{minipage}
  \begin{minipage}{5.5cm}
    \vspace{1.0cm}
    \input{graf07.eepic}
  \end{minipage}

  A graph paper can be added simply by calling the member function
  \textit{grid()}.

  \begin{minipage}{7.5cm}
    \begin{lstlisting}
  #include <epix.h>
  int main()
  {
    CartesianCoord cs(-2,4,-1,5);
    picture(P(12,12));
    unitlength("5mm");
    begin();

    cs.grid();
    cs.draw();

    end();
    return 0;
  }
    \end{lstlisting}
  \end{minipage}
  \begin{minipage}{5.5cm}
    \input{graf08.eepic}
  \end{minipage}


  The position of the paper is set by calling the member function
  \textit{gridPosition()} with a bold line crossing point as parameter
  (default: axis crossing point). The density of the gridlines is set
  by calling the member function \textit{gridDensity()} (see reference
  section for more details):

  \begin{minipage}{7.5cm}
    \begin{lstlisting}
  #include <epix.h>
  static char* ylabels[] = {"10",NULL};
  int main()
  {
    CartesianCoord cs(-2,4,-10,50);
    picture(P(12,12));
    unitlength("5mm");
    begin();

    cs.yMarks(10,1);
    cs.yLabels(ylabels);
    cs.grid();
    cs.gridDensity(1,9,5,10,9,0);
    cs.draw();

    end();           
    return 0;
  }
    \end{lstlisting}
  \end{minipage}
  \begin{minipage}{5.5cm}
    \vspace{1.0cm}
    \input{graf09.eepic}
  \end{minipage}


  \newpage
  \section{Logarithmic Coordinate Systems}

  Logarithmic coordinate systems are made in much the same way as the
  cartesian coordinate system.

  \subsection{SingleLog Cordinate System}

  A single logarithmic coordinate system has a horizontal linear axis
  and a vertical logarithmic axis.

  The declaration statement

  \begin{lstlisting}
    SingleLogCoord cs(0,10,1,100);
  \end{lstlisting}

  creates a logarithmic coordinate system with x-bounds $0\ldots10$
  and y-bounds $1\ldots10$. To draw the coordinate system using the
  standard layout set by the constructor of \textit{SingleLogCoord},
  execute the member function \textit{draw} in \textit{cs}:

  \begin{lstlisting}
    cs.draw();
  \end{lstlisting}

  Here is an example of a complete C++ program making such a standard
  coordinate system:

  \begin{minipage}{7.5cm}
    \begin{lstlisting}
  #include "epix.h"

  int main()
  {
    SingleLogCoord sc(0,8,1,100);

    picture(P(12,12));
    unitlength("5mm");
    begin();

    sc.draw();

    end();
    return 0;
  }
    \end{lstlisting}
  \end{minipage}
  \begin{minipage}{5.5cm}
    \input{graf10.eepic}
  \end{minipage}

  The logaritmic axis need not contain only full decades. The
  following C++ program is used to create a coordinate system with one
  and a "half" decades from 0.1 to 5:

  \begin{minipage}{7.5cm}
    \begin{lstlisting}
  #include "epix.h"

  int main()
  {
    SingleLogCoord sc(0,8,0.1,5);
    picture(P(12,12));
    unitlength("5mm");
    begin();

    sc.logStyle(3);
    sc.draw();

    end();
    return 0;
  }
    \end{lstlisting}
  \end{minipage}
  \begin{minipage}{5.5cm}
    \input{graf11.eepic}
  \end{minipage}


  In this case the standard logaritmic style number 3 is used to set
  the division of the logaritmic axis.

  There are six standard styles giving an increasing number of
  division points:

  \begin{center}
    \input{stylegraf1.eepic}
  \end{center}
  \begin{center}
    \input{stylegraf2.eepic}
  \end{center}
  \begin{center}
    \input{stylegraf3.eepic}
  \end{center}
  \begin{center}
    \input{stylegraf4.eepic}
  \end{center}
  \begin{center}
    \input{stylegraf5.eepic}
  \end{center}
  \begin{center}
    \input{stylegraf6.eepic}
  \end{center}

  The default \textit{stdstyle} is 4.

  The start of each decade is labeled as shown inside the range 0.01
  to 1000.  Outside this range powers of ten are used.

  To customize marking and labelling call the member function
  \textit{logStyle()} as shown here:

  \begin{minipage}{8.0cm}
    \begin{lstlisting}
  #include "epix.h"

  Mark mrk[] =
   {
    {1,"1000",3},{2,"2000",1},
    {3,"3000",1},{4,"4000",1},
    {5,"5000",0},{6,"6000",0},
    {7,"7000",0},{8,"8000",0},
    {9,"9000",0},{1,"10000",3},
    {2,"20000",1},{3,"30000",1},
    {4,"40000",1},{5,"50000",0},
    {6,"60000",0},{7,"70000",0},
    {8,"80000",0},{9,"90000",0},
    {1,"100000",0},{10,NULL,0}
   };
  int main()
  {
    SingleLogCoord
      cs(-1,10,1000,100000);
    picture(P(12,24));
    unitlength("5mm");
    begin();

    cs.logStyle(mrk);
    cs.draw();

    end();
    return 0;
  }
    \end{lstlisting}
  \end{minipage}
  \begin{minipage}{5.5cm}
    \input{graf12.eepic}
  \end{minipage}

  The variable \textit{mrk} is a table of structures each containing
  three fields. The first field holds the position of the mark (the
  position within the actual decade). The next field is the label to
  be printed. And the third is the number of subdivisons following
  this mark. The table ends with position 10.

  \newpage
  Another example:

  \begin{minipage}{8.0cm}
    \begin{lstlisting}
  #include "epix.h"

  Mark mrk[] =
   {
    {1,NULL,9},
    {2,"\\scriptsize 2",9},
    {3,"\\scriptsize 3",9},
    {4,"\\scriptsize 4",9},
    {5,"\\scriptsize 5",4},
    {6,"\\scriptsize 6",4},
    {7,"\\scriptsize 7",4},
    {8,"\\scriptsize 8",4},
    {9,"\\scriptsize 9",4},
    {10,NULL,0}
   };

  int main()
  {
    SingleLogCoord cs(0,10,1,100);
    picture(P(12,24));
    unitlength("5mm");
    begin();

    cs.logStyle(mrk);
    cs.labelAttr("");
    cs.draw();

    end();
    return 0;
  }
    \end{lstlisting}
  \end{minipage}
  \begin{minipage}{5.5cm}
    \input{graf13.eepic}
  \end{minipage}

  \textit{mrk} is here a one-decade table. It is used for both decades
  on the axis. The NULL label at position 1 has the effect that the
  current decade value is printed (here 1.0, 10 and 100).

  The member function \textit{labelAttr()} can be used to set any
  attribute, for instance $\backslash$it and/or $\backslash$tiny. But
  note that you have to use a double backslash in a C++ string,
  because in C++ strings (like in \TeX) a backslash is a command
  introducer. In the example the attribute is set to the empty string
  so that the numbers at the start of each decade are printed in
  normal size (default attribute string is
  "$\backslash\backslash$scriptsize").  All other numbers are printed
  using $\backslash$scriptsize as specified in the Mark table.

  \newpage
  A graph paper can be added in the same way as for the cartesian
  coordinate system:

  \begin{minipage}{7.5cm}
    \begin{lstlisting}
  #include "epix.h"
  int main()
  {
    SingleLogCoord cs(0,10,1,10);
    picture(P(12,12));
    unitlength("5mm");
    begin();

    cs.logStyle(6);
    cs.grid();
    cs.draw();

    end();
    return 0;
  }
    \end{lstlisting}
  \end{minipage}
  \begin{minipage}{5.5cm}
    \input{graf14.eepic}
  \end{minipage}

  \subsection{DoubleLog Cordinate System}

  A double logarithmic coordinate system has two logarithmic
  axisis. It is handled in much the same way as the Single Logarithmic
  Coordinate System.  Just change \textit{Single} to \textit{Double}.

  Here is an example:

  \begin{minipage}{7.7cm}
    \begin{lstlisting}
  #include "epix.h"
  int main()
  {
    DoubleLogCoord cs(0.1,10,1,100);

    picture(P(12,12));
    unitlength("5mm");
    begin();

    cs.logStyle(4);
    cs.grid();
    cs.draw();

    end();
    return 0;
  }
    \end{lstlisting}
  \end{minipage}
  \begin{minipage}{5.5cm}
    \input{graf15.eepic}
  \end{minipage}

  \subsection{Graphs in Logarithmic Cordinate System}

  In the coordinate systems the member functions \textit{plot()},
  \textit{adplot()}, \textit{label()} and \textit{marker()} are used
  in the same way as the corresponding ePiX functions. But they can be
  used in logarithmic coordinate systems, too.


  %Here is a simple example of the use of these:

  \begin{minipage}{7.5cm}
    \begin{lstlisting}
  #include "epix.h"
  double f(double x)
  { return 1.5*pow(1.2,x); }
  int main()
  {
    SingleLogCoord cs(0,10,1,10);
    picture(P(12,12));
    unitlength("5mm");
    begin();
    cs.logStyle(6);
    cs.grid();
    cs.draw();
    cs.plot((*f),0,10,10);
    cs.label(P(5,f(5)),P(0,0),
       "$1.5\\cdot1.2^x$",br);
    end();
    return 0;
  }
    \end{lstlisting}
  \end{minipage}
  \begin{minipage}{5.5cm}
    \input{graf16.eepic}
  \end{minipage}

  By changing \textit{SingleLog} to \textit{Cartesian} (and doing some
  other natural changes), you get:

  \begin{minipage}{7.5cm}
    \begin{lstlisting}
  #include "epix.h"
  double f(double x)
  { return 1.5*pow(1.2,x); }
  int main()
  {
    CartesianCoord cs(0,10,0,10);
    picture(P(12,12));
    unitlength("5mm");
    begin();
    cs.grid();
    cs.draw();
    cs.plot((*f),0,10,10);
    cs.label(P(5,f(5)),P(0,0),
       "$1.5\\cdot1.2^x$",br);
    end();
    return 0;
  }
    \end{lstlisting}
  \end{minipage}
  \begin{minipage}{5.5cm}
    \input{graf17.eepic}
  \end{minipage}


  \subsection{Printing Coordinate Axis}

  Coordinate axis can be printed alone. The output of the program
  below shows the theory behind logarithmic axis:

  \begin{lstlisting}
    #include "epix.h"
    static char* linLabels[] =
      {"\\scriptsize 0","\\scriptsize 0.1","\\scriptsize 0.2",
       "\\scriptsize 0.3","\\scriptsize 0.4","\\scriptsize 0.5",
       "\\scriptsize 0.6","\\scriptsize 0.7","\\scriptsize 0.8",
       "\\scriptsize 0.9","\\scriptsize 1.0",NULL
      };
    int main()
    {
      HorizLinAxis linAs;
      HorizLogAxis logAs;

      bounding_box(P(0,0),P(1,1.17));
      picture(P(10,1.5));
      unitlength("1cm");
      begin();

      linAs.nummrk = 11;
      linAs.distmrk = 0.1;
      linAs.mrk1 = 0; linAs.mrkpos = POSITIVE;
      linAs.labels = linLabels; linAs.labpos = POSITIVE;
      linAs.name = "$\\log(x)$"; linAs.nampos = POSITIVE;
      linAs.draw(P(0,0.75),1.17);

      logAs.drawaxis = false;
      logAs.arrow = true; logAs.arrowextra = 5;
      logAs.stdstyle = 6;
      logAs.mrkpos = NEGATIVE;
      logAs.name = "$x$";
      logAs.draw(P(0,0.75),1);

      end();
      return 0;
    }
  \end{lstlisting}

  \begin{minipage}{5.5cm}
    \input{graf18.eepic}
  \end{minipage}

  See \textit{epix\_ext.h} for a description of the fields in the axis
  structures.

  %\newpage
  \section{Hatched figures}

  \subsection{Hatched area}

  Commands to draw a hatched area between two graphs (or one graph and
  the x-axis) in a cartesian coordinate system are generated by
  calling the member functions:

  \begin{lstlisting}
  void hatchArea(double angle,double dist,
                   double f(double),double g(double),
                   double a, double b, int n);

  void hatchArea(double angle,double dist,
                   double f(double),double a,double b, int n)
  \end{lstlisting}

  The first two parameters \textit{angle} and \textit{dist} specifies
  the angle of and the distance between the hatch lines. The next are
  the function(s), the start and the end point and the number of
  points used to draw the graphs.

  \begin{minipage}{7.5cm}
    \begin{lstlisting}
  #include "epix.h"

  double f(double x)
  {
    return 4.5-(x-3)*(x-3)/2;
  }

  int main()
  {
    CartesianCoord cs(-1,7,-2,6);
    picture(P(12,12));
    unitlength("5mm");
    begin();

    cs.draw();
    cs.hatchArea(150,0.2,f,std::sin,0,5.5,100);

    end();
    return 0;
  }
    \end{lstlisting}
  \end{minipage}
  \begin{minipage}{5.5cm}
    \input{graf20.eepic}
    \vspace{1.8cm}
  \end{minipage}

  %\pagebreak

  \subsection{Hatched polygon}

  The member function \textit{hatchArea()} is based on routines for drawing
  hatched polygons. They are generated by the procedures:

  \begin{lstlisting}
    void hatch_polygon(double angle,double dist,
                       int corners,triple* points);
    void hatch_polygon(double angle,double dist,
                       int corners,triple& points, ...);

    void hatch(double angle,double dist,
               int corners,triple* points);
    void hatch(double angle,double dist,
               int corners,triple& points, ...);
  \end{lstlisting}

  The difference between \textit{hatch\_polygon()} and
  \textit{hatch()} is that \textit{hatch()} will not draw the
  polygon. The meaning of the first two parameters are the same as for
  the \textit{hatchArea()} member function. The next parameter is the
  number of corners in the polygon and then the corners.

  \begin{minipage}{7.5cm}
    \begin{lstlisting}
  #include "epix.h"

  triple corners[] =
    {P(-2,-3),P(-3,1),P(1,0),P(2,3),P(4,-2),P(-1,-1)};

  int main()
  {
    bounding_box(P(-4,-3),P(4,2));
    picture(P(12,7.5));
    unitlength("5mm");

    begin();

    hatch_polygon(45,0.2,6,corners);

    end();

    return 0;
  }
    \end{lstlisting}
  \end{minipage}
  \begin{minipage}{5.5cm}
    \vspace{1.5cm}
    \input{graf22.eepic}
  \end{minipage}

  As is seen from the example, the polygon need not be convex. But the
  result is unpredictable if some of the edges cosses each other.

  \section{Setting Line Thickness}

  The line thickness used to draw axis (bold and normal), bold, medium
  and fine grid lines can be changed simply by setting variables to
  new values.  These variables and their standard values are:

  \begin{lstlisting}
    char* normalThickness     = "\\thinlines";
    char* normalAxisThickness = "\\thinlines";
    char* boldAxisThickness   = "\\thicklines";
    char* boldGridThickness   = "\\thinlines";
    char* mediumGridThickness = "\\allinethickness{0.075mm}";
    char* fineGridThickness   = "\\allinethickness{0.03mm}";
    char* crossHatchThickness = "\\allinethickness{0.04mm}";
  \end{lstlisting}


  \section{ePiX\_ext reference}

  \subsection{Member functions for all coordinate systems}

  \begin{MemberDescr}
    void bounds(double xmin,double xmax,double ymin,double ymax)\\
    \memdsc
      Set bounds of coordinate system.\\
    \meminp
      xmin: lower x-bound of picture.\\
      xmax: upper x-bound of picture.\\
      ymin: lower y-bound of picture.\\
      ymax: upper y-bound of picture.\\
  \end{MemberDescr}

  \begin{MemberDescr}
    void draw()\\
    \memdsc
      Draw the coordinate system.\\
    \meminp
      none.\\
  \end{MemberDescr}

  \begin{MemberDescr}
    void showAxis(bool b)\\
    void showAxis(bool bx,bool by)\\
    \memdsc
      Draw axis or not.\\
    \meminp
      b:  if true draw both x- and y-axis.\\
      bx: if true draw x-axis.\\
      by: if true draw y-axis.\\
    \memdef
      \textbf{true} for both axis.\\
  \end{MemberDescr}

  \begin{MemberDescr}
     void showArrow(bool b)\\
     void showArrow(bool xb, bool yb)\\
    \memdsc
      Draw axis with or without an arrow.\\
    \meminp
      b:  if true draw arrow on both x- and y-axis.\\
      bx: if true draw arrow on x-axis.\\
      by: if true draw arrow on y-axis.\\
    \memdef
      \textbf{true} for both axis.\\
  \end{MemberDescr}

  \begin{MemberDescr}
     void xName(Position pos)\\
     void xName(char\pntr name,Position pos)\\
    \memdsc
      Place name at the end of the x-axis.\\
    \meminp
      name: name of the x-axis.\\
      pos:  position of name (POSITIVE, NEGATIVE or CENTER).\\
    \memdef
      name: "(1)" for linear axis and \textbf{NULL} for logarithmic.\\
      pos:  NEGATIVE.\\
  \end{MemberDescr}

  \begin{MemberDescr}
     void yName(Position pos)\\
     void yName(char\pntr name,Position pos)\\
    \memdsc
      Place name at the end of the y-axis.\\
    \meminp
      name: name of the y-axis.\\
      pos:  position of name (POSITIVE, NEGATIVE or CENTER).\\
    \memdef
      name: "(2)" for linear axis and \textbf{NULL} for logarithmic.\\
      pos:  NEGATIVE.\\
  \end{MemberDescr}

  \begin{MemberDescr}
     void labelPos(Position pos)\\
     void labelPos(Position xpos,Position ypos)\\
    \memdsc
      Set position of (mark) labels on the axis.\\
    \meminp
      pos:  position for labels on both x- and y-axis.\\
      xpos: position for labels on x-axis.\\
      ypos: position for labels on y-axis.\\
    \memdef
      NEGATIVE for both axis.\\
  \end{MemberDescr}

  \begin{MemberDescr}
    void markLayout(double length,Position pos=CENTER)\\
    void markLayout(double lx,Position px,double ly,Position py)\\
    \memdsc
      Set length and position of marks on the axis.\\
    \meminp
      length: length (in pt) of marks on both x-axis and y-axis.\\
      pos:    position of marks on both x-axis and y-axis.\\
      lx:     length (in pt) of marks on x-axis.\\
      px:     position of marks on x-axis.\\
      ly:     length (in pt) of marks on y-axis.\\
      py:     position of marks on y-axis.\\
    \memdef
      length: 3 pt.\\
      pos:    CENTER.\\
  \end{MemberDescr}

  \begin{MemberDescr}
    void boldAxis(bool b=true)\\
    void boldAxis(bool bx,bool by)\\
    \memdsc
      draw axis with bold line.\\
    \meminp
      b:  if true draw both x- and y-axis with bold line.\\
      bx: if true draw x-axis with bold line.\\
      by: if true draw y-axis with bold line.\\
    \memdef
      For both axis \textbf{true} if grid is selected, otherwise \textbf{false}.\\
  \end{MemberDescr}

  \begin{MemberDescr}
    void adplot(double f(double), double xmin, double xmax, int n)\\
    void plot(double f(double), double xmin, double xmax, int n)\\
    \memdsc
      Works like ePiX adplot and plot but takes care of the the actual type of
      coordinate system (cartesian or logarithmic).\\
    \meminp
      f: function to draw graph of.\\
      xmin: lower x-bound of graph.\\
      xmax: upper x-bound of graph.\\
      n: number of points used to draw (straight lines between each point).\\
  \end{MemberDescr}

  \begin{MemberDescr}
    void label(triple base,triple offset,char\pntr text)\\
    void label(triple base,char *text)\\
    void label(triple base,triple offset,char\pntr text,epix\_label\_posn pos)\\
    \memdsc
      Works like ePiX label but takes care of the the actual type of
      coordinate system (cartesian or logarithmic).\\
    \meminp
      base:   base point of label.\\
      offset: offset (in pt) of label.\\
      text:   text to write.\\
      pos:    position relative to base (c,r,tr,rt,t,tl,lt,l,bl,lb,b,br,rb).\\
  \end{MemberDescr}


  \subsection{Member functions for class CartesianCoord}

  \begin{MemberDescr}
    CartesianCoord()\\
    CartesianCoord(double xmin,double xmax,double ymin,double ymax)\\
    \memdsc
      Constructor for class CartesianCoord.\\
    \meminp
      xmin: lower x-bound of picture (default -1).\\
      xmax: upper x-bound of picture (default 10).\\
      ymin: lower y-bound of picture (default -1).\\
      ymax: upper y-bound of picture (default 10).\\
  \end{MemberDescr}

  \begin{MemberDescr}
    void axisCross(triple point)\\
    \memdsc
      Set the axis crossing point.\\
    \meminp
      point: crossing point for the axis.\\
    \memdef
      (0,0).\\
  \end{MemberDescr}

  \begin{MemberDescr}
    void axisBounds(triple upperRight)\\
    void axisBounds(triple lowerLeft,triple upperRight)\\
    \memdsc
      Set bouds of axis in coordinate system.\\
    \meminp
      lowerLeft:  lower left corner of axis.\\
      upperRight: upper right corner of axis.\\
    \memdef
      Bounds of coordinate system.\\
  \end{MemberDescr}

  \begin{MemberDescr}
    void xMarks(double first,int num)\\
    void xMarks(double first,int num,double dist)\\
    \memdsc
      Specify number and position of marks on x-axis.\\
    \meminp
      first: position of first mark.\\
      num:   number of marks to set.\\
      dist:  distance between marks (default is 1.0).\\
    \memdef
      first: 1.\\
      num:   1.\\
      dist:  1.\\
  \end{MemberDescr}

  \begin{MemberDescr}
    void yMarks(double first,int num).\\
    void yMarks(double first,int num,double dist).\\
    \memdsc
      Specify number and position of marks on y-axis.\\
    \meminp
      first: position of first mark.\\
      num:   number of marks to set.\\
      dist:  distance between marks (default is 1.0).\\
    \memdef
      first: 1.\\
      num:   1.\\
      dist:  1.\\
  \end{MemberDescr}

  \begin{MemberDescr}
    void xLabels(char\ppntr lab)\\
    void xLabels(char\ppntr lab,double first)\\
    void xLabels(char\ppntr lab,double first,double dist)\\
    \memdsc
      Labels to place on x-axis.\\
    \meminp
      lab:   the labels in a \textbf{NULL}-terminated array of string pointers.\\
      first: position of label (default is position of first mark).\\
      dist:  distance between labels (default distance between marks).\\
    \memdef
      lab: \{"1",\textbf{NULL}\}.\\
      first: first mark.\\
      dist: distance between marks.\\
  \end{MemberDescr}

  \begin{MemberDescr}
    void yLabels(char\ppntr lab)\\
    void yLabels(char\ppntr lab,double first)\\
    void yLabels(char\ppntr lab,double first,double dist)\\
    \memdsc
      Labels to place on y-axis.\\
    \meminp
      lab:   the labels in a \textbf{NULL}-terminated array of string pointers.\\
      first: position of label (default is position of first mark).\\
      dist:  distance between labels (default distance between marks).\\
    \memdef
      lab: \{"1",\textbf{NULL}\}.\\
      first: first mark.\\
      dist: distance between marks.\\
  \end{MemberDescr}

  \begin{MemberDescr}
    void xBroken()\\
    \memdsc
      Make x-axis broken (to show that axis cross is not zero).\\
    \meminp
      None.\\
  \end{MemberDescr}

  \begin{MemberDescr}
    void yBroken()\\
    \memdsc
      Make y-axis broken (to show that axis cross is not zero).\\
    \meminp
      None.\\
  \end{MemberDescr}

  \begin{MemberDescr}
    void grid()\\
    void grid(double xmin,double xmax,double ymin,double ymax)\\
    \memdsc
      Add a grid to the coordinate system. Draw coordinate axis in bold unless
      otherwise specified.\\
    \meminp
      xmin: left edge of grid (default is left edge of coordinate system).\\
      xmax: right edge of grid (default is right edge of coordinate system).\\
      ymin: bottom edge of grid (default is bottom edge of coordinate system).\\
      ymax: top edge of grid (default is top edge of coordinate system).\\
  \end{MemberDescr}

  \begin{MemberDescr}
    void gridPosition(triple boldCross)\\
    \memdsc
      Set position of grid in coordinate system.\\
    \meminp
      boldCross: a crossing point for bold lines.\\
    \memdef
      Axis crossing point.\\
  \end{MemberDescr}

  \begin{MemberDescr}
    void gridDensity(double dist)\\
    void gridDensity(double dist,int num)\\
    void gridDensity(double dist,int num,int mid)\\
    void gridDensity(double xd,int xn,int xm,double yd,int yn,int ym)\\
    \memdsc
      Set density of grid lines.\\
    \meminp
      dist: distance between grid lines.\\
      num:  number of fine lines (default is 9).\\
      mid:  the fine line to draw half bold (default is 0 meaning no half bold).\\
      xd,xn,xm: dist, num and mid for x-axis.\\
      yd,yn,ym: dist, num and mid for y-axis.\\
    \memdef
      dist: 1.\\
      num:  9.\\
      mid:  5.\\
  \end{MemberDescr}

  \begin{MemberDescr}
    void hatchArea(double v,double d,double f(double),\\
    \phantom{void hatchArea(}double a,double b,int n)\\
    \memdsc
      Hatched area between graph and x-axis.\\
    \meminp
      v: angle between x-axis and hatch line.\\
      d: distance between hatch lines.\\
      f: graph function.\\
      a: left edge af area.\\
      b: right edge of area.\\
      n: number of points used to draw the graphs (see plot and adplot).\\
  \end{MemberDescr}

  \begin{MemberDescr}
    void hatchArea(double v,double d,double f(double),double g(double)\\
    \phantom{void hatchArea(}double a,double b,int n)\\
    \memdsc
      Hatched area between two graphs.\\
    \meminp
      v: angle between x-axis and hatch line.\\
      d: distance between hatch lines.\\
      f: upper graph function.\\
      g: lower graph function.\\
      a: left edge af area.\\
      b: right edge of area.\\
      n: number of points used to draw the graphs (see plot and adplot).\\
  \end{MemberDescr}

  \subsection{Member functions for class SingleLogCoord}

  \begin{MemberDescr}
    SingleLogCoord()\\
    SingleLogCoord(double xmin,double xmax,double ymin,double ymax)\\
    \memdsc
      Constructor for class SingleLogCoord.\\
    \meminp
      xmin: lower x-bound of picture (default -1).\\
      xmax: upper x-bound of picture (default 10).\\
      ymin: lower y-bound of picture (default 1).\\
      ymax: upper y-bound of picture (default 100).\\
  \end{MemberDescr}

  \begin{MemberDescr}
    void axisCross(triple point)\\
    \memdsc
      Set crossing point for coordinate axis. Set to (xmin,ymin) by
      the constructor.\\
    \meminp
      point: crossing point.\\
    \memdef
      (xmin,ymin).\\
  \end{MemberDescr}

  \begin{MemberDescr}
    void axisBounds(triple upperRight)\\
    void axisBounds(triple lowerLeft,triple upperRight)\\
    \memdsc
      Set bounds of axis in coordinate system. Set by constructor to bounds
      of picture.\\
    \meminp
      lowerLeft:  lower left corner of axis.\\
      upperRight: upper right corner of axis.\\
    \memdef
      lowerLift: (xmin,ymin).\\
      upperRight: (xmax,ymax).\\
  \end{MemberDescr}

  \begin{MemberDescr}
    void xMarks(double first,int num)\\
    void xMarks(double first,int num,double dist)\\
    \memdsc
      Specify number and position of marks on x-axis.\\
    \meminp
      first: position of first mark.\\
      num:   number of marks to set.\\
      dist:  distance between marks (default is 1.0).\\
    \memdef
      first: 1.\\
      num:   1.\\
      dist:  1.\\
  \end{MemberDescr}

  \begin{MemberDescr}
    void yMarks(double first,double last)\\
    \memdsc
      Set bounds of marks on logarithmic axis.\\
    \meminp
      first: lower bound of marks.\\
      last: upper bounds marks.\\
    \memdef
      first: ymin.\\
      last: ymax.\\
  \end{MemberDescr}

  \begin{MemberDescr}
    void xLabels(char\ppntr lab)\\
    void xLabels(char\ppntr lab,double first)\\
    void xLabels(char\ppntr lab,double first,double dist)\\
    \memdsc
      Labels to place on x-axis.\\
    \meminp
      lab:   the labels in a \textbf{NULL}-terminated array of string pointers.\\
      first: position of label (default is position of first mark).\\
      dist:  distance between labels (default distance between marks).\\
    \memdef
      lab: \{"1",\textbf{NULL}\}.\\
      first: first mark.\\
      dist: distance between marks.\\
  \end{MemberDescr}

  \begin{MemberDescr}
    void showArrow(bool b)\\
    void showArrow(double extra)\\
    void showArrow(bool b,double extra)\\
    \memdsc
      Draw axis with or without an arrow and add extra length to logarithmic
      axis.\\

      \textbf{void showArrow(bool b)} will only affect x-axis.\\
      \textbf{void showArrow(double extra)} will only affect y-axis.\\
      \textbf{void showArrow(bool b,double extra)} affects both axis.\\
    \meminp
      b:  if true draw arrow on x-axis.\\
      extra: add extra length (in pt) to y-axis and draw arrow.\\
    \memdef
      Arrow on x-axis, no arrow on y-axis.\\
  \end{MemberDescr}

  \begin{MemberDescr}
    void xBroken()\\
    \memdsc
      Make x-axis broken (to show that axis cross is not zero).\\
    \meminp
      None.\\
    \memdef
      No broken axis.\\
  \end{MemberDescr}


  \begin{MemberDescr}
    void logStyle(int style)\\
    void logStyle(Mark\pntr mrk)\\
    \memdsc
      Set style of logarithmic axis. \textit{Mark} is defined as:\\
    \vspace{-1pc}
\begin{lstlisting}
  class Mark
  {
    public:
     double position; // position of main mark
     char*  label;    // label to print at main mark
     int    numsub;   // number of submarks until next main mark
  };

\end{lstlisting}

    \meminp
      style: standard style number ($1\ldots6$).\\
      mrk:   pointer to an array of marks describing the logarithmic axis.\\
    \memdef
      Standard style 4.\\
  \end{MemberDescr}

  \begin{MemberDescr}
    void labelAttr(char\pntr attr)\\
    \memdsc
      Label attribute to use if label is not specified (\textbf{NULL}).\\
    \meminp
      attr: attribute string.\\
    \memdef
      attr = "$\backslash\backslash$scriptsize".\\
  \end{MemberDescr}

  \begin{MemberDescr}
    void grid()\\
    void grid(double xmin,double xmax)\\
    \memdsc
      Add a grid to the coordinate system. Draw coordinate axis in bold unless
      otherwise specified.\\
    \meminp
      xmin: left edge of grid (default is left edge of coordinate system).\\
      xmax: right edge of grid (default is right edge of coordinate system).\\
    \memdef
      xmin: left edge of coordinate system.\\
      xmax: right edge of coordinate system.\\
  \end{MemberDescr}

  \begin{MemberDescr}
    void gridPosition(double crossX)\\
    \memdsc
      Set position of grid lines perpendicular to x-axis.\\
    \meminp
      crossX: position for a bold line.\\
    \memdef
      logarithmic axis position.\\
  \end{MemberDescr}

  \begin{MemberDescr}
    void gridDensity(double dist)\\
    void gridDensity(double dist,int num)\\
    void gridDensity(double dist,int num,int mid)\\
    \memdsc
      Density of grid lines perpendicular to x-axis.\\
    \meminp
      dist: distance between grid lines.\\
      num:  number of fine lines (default is 9).\\
      mid:  the fine line to draw half bold (default 0 meaning no half bold).\\
    \memdef
      dist: 1.\\
      num:  9.\\
      mid:  5.\\
  \end{MemberDescr}


  \subsection{Member functions for class DoubleLogCoord}

  \begin{MemberDescr}
    DoubleLogCoord()\\
    DoubleLogCoord(double xmin,double xmax,double ymin,double ymax)\\
    \memdsc
      Constructor for class DoubleLogCoord.\\
    \meminp
      xmin: lower x-bound of picture (default 1).\\
      xmax: upper x-bound of picture (default 100).\\
      ymin: lower y-bound of picture (default 1).\\
      ymax: upper y-bound of picture (default 100).\\
  \end{MemberDescr}

  \begin{MemberDescr}
    void axisCross(triple point)\\
    \memdsc
      Set crossing point for coordinate axis.\\
    \meminp
      point: crossing point.\\
    \memdef
      (xmin,ymin).\\
  \end{MemberDescr}

  \begin{MemberDescr}
    void axisBounds(triple upperRight)\\
    void axisBounds(triple lowerLeft,triple upperRight)\\
    \memdsc
      Set bounds of axis in coordinate system. Set by constructor to bounds
      of picture.\\
    \meminp
      lowerLeft:  lower left corner of axis.\\
      upperRight: upper right corner of axis.\\
    \memdef
      lowerLeft: (xmin,ymin).\\
      upperRight: (xmax,ymax).\\
  \end{MemberDescr}

  \begin{MemberDescr}
    void xMarks(double first,double last)\\
    \memdsc
      Set bounds of marks on horizontal logarithmic axis.\\
    \meminp
      first: lower bound of marks.\\
      last: upper bounds marks.\\
    \memdef
      first: xmin.\\
      last: xmax.\\
  \end{MemberDescr}

  \begin{MemberDescr}
    void yMarks(double first,double last)\\
    \memdsc
      Set bounds of marks on vertical logarithmic axis.\\
    \meminp
      first: lower bound of marks.\\
      last: upper bounds marks.\\
    \memdef
      first: ymin.\\
      last: ymax.\\
  \end{MemberDescr}

  \begin{MemberDescr}
    void showArrow(double extra)\\
    void showArrow(bool bx,double extrax,bool by,double extray)\\
    \memdsc
      Draw axis with or without an arrow and add extra length to logarithmic
      axis.\\
    \meminp
      extra: add extra length (in pt) and draw arrow (for both axis).\\
      extrax: add extra length (in pt) to horizontal axis and draw arrow.\\
      extray: add extra length (in pt) to vertical axis and draw arrow.\\
    \memdef
      No arrows.\\
  \end{MemberDescr}

  \begin{MemberDescr}
    void logStyle(int style)\\
    void logStyle(Mark\pntr mrk)\\
    void logStyle(int stylex,int styley)\\
    void logStyle(Mark\pntr mrkx,Mark\pntr mrky)\\
    void logStyle(int stylex,Mark\pntr mrky)\\
    void logStyle(Mark\pntr mrkx,int styley)\\
    \memdsc
      Set style of axis. \textit{Mark} is defined as:\\
    \vspace{-1pc}
\begin{lstlisting}
  class Mark
  {
    public:
     double position; // position of main mark
     char*  label;    // label to print at main mark
     int    numsub;   // number of submarks until next main mark
  };

\end{lstlisting}

    \meminp
      style: standard style ($1\ldots6$) to be used for both axis.\\
      stylex: standard style ($1\ldots6$) to be used for horizontal axis.\\
      styley: standard style ($1\ldots6$) to be used for vertical axis.\\
      mrk:   pointer to an array of marks describinging both axis.\\
      mrkx:   pointer to an array of marks describinging horizontal axis.\\
      mrky:   pointer to an array of marks describinging vertical axis.\\
    \memdef
      Standard style 4 for both axis.\\
  \end{MemberDescr}

  \begin{MemberDescr}
    void labelAttr(char\pntr attr)\\
    void labelAttr(char\pntr attrx,char\pntr attry)\\
    \memdsc
      Label attribute to use if label is not specified (\textbf{NULL}).\\
    \meminp
      attr: attribute string for both axis.\\
      attrx: attribute string for horizontal axis.\\
      attry: attribute string for vertical axis.\\
    \memdef
      attr = "$\backslash\backslash$scriptsize".\\
  \end{MemberDescr}

  \begin{MemberDescr}
    void grid()\\
    \memdsc
      Add a grid to the coordinate system. Draw coordinate axis in bold unless
      otherwise specified.\\
    \meminp
      None.\\
  \end{MemberDescr}

\end{document}

